% Rijen zonder limieten (pakket 1, 2 en 3). Monet: een rij populieren langs het water, stam na stam.
% Motief: de trappen uit het probleem van hoofdstuk 1. Trap n is trap n - 1 met een kolom van n blokjes
% erbij (fel): 1, 3, 6, 10, 15 blokjes. Wat erbij komt, 1, 2, 3, 4, 5, is een rekenkundige rij.
\def\omslagbestand{omslag-populieren.jpg}
\def\omslagbron{Claude Monet, \textit{Populieren aan de Epte}, 1891 (publiek domein)}
\def\omslagrand{10.5 + 0.5*\x/21}
\def\omslagmotief{%
  \randlijn
  \def\blok{.66}%
  \foreach \n/\links/\aantal in {1/8.9/1,2/9.98/3,3/11.72/6,4/14.12/10,5/17.18/15}{%
    \foreach \k in {1,...,\n}{%
      \foreach \h in {1,...,\k}{%
        \ifnum\k=\n
          \fill[fel,draw=accent,line width=.6pt] ({\links+(\k-1)*\blok},{5+(\h-1)*\blok}) rectangle ++(\blok,\blok);
        \else
          \fill[white,fill opacity=.2] ({\links+(\k-1)*\blok},{5+(\h-1)*\blok}) rectangle ++(\blok,\blok);
          \draw[white,line width=.6pt] ({\links+(\k-1)*\blok},{5+(\h-1)*\blok}) rectangle ++(\blok,\blok);
        \fi
      }%
    }%
    \node[text=white,font=\small,inner sep=1pt,anchor=north] at ({\links+\n*\blok/2},4.85) {$\aantal$};
  }%
}
